\documentclass[10pt,a4paper]{amsart}
\usepackage[margin=19mm]{geometry}
\usepackage{amsmath,amssymb,mathtools}
\usepackage[scr=rsfs,cal=euler]{mathalfa}
\usepackage{xfrac}
\usepackage[unicode,hidelinks,bookmarks=false]{hyperref}
\newcommand{\ie}{i.e.\ }
\newcommand{\cf}{cf.\ }
\newcommand{\resp}{respectively\ }
\newcommand{\Subsection}{\subsection}
\providecommand{\nobreakdash}{\nobreak\hbox{-}}
\setlength{\emergencystretch}{2em}
\multlinegap0pt
\usepackage{mathtools}
\usepackage{enumerate}
\usepackage{caption}
\newtheorem{proposition}{Proposition}
\newcommand{\cal}{\mathcal}
\newcommand{\clb}{{\cal B}}
\newcommand{\clh}{{\cal H}}
\newcommand{\raro}{\rightarrow}
\title{Isometric actions of compact quantum groups on graph $C^{\ast}$-algebras}
\author{Soumalya Joardar, Jitender Sharma}
\date{Source: arXiv:2409.19376v1 (CC BY 4.0)}
\begin{document}
\maketitle

\noindent\emph{Source and licence.} Isometric actions of compact quantum groups on graph $C^{\ast}$-algebras. Version arXiv:2409.19376v1 (CC BY 4.0), distributed by arXiv under the Creative Commons Attribution 4.0 International licence, http://creativecommons.org/licenses/by/4.0/, as recorded in the arXiv licence metadata for this exact version and shown on the abstract page https://arxiv.org/abs/2409.19376v1. Full licence notice: \texttt{licenses/arxiv-2409.19376-CC-BY-4.0.txt}.\par\smallskip

\noindent\textit{Editorial scope.} Counted unit: the article's proposition implementation (source0.tex lines 792--795). The article's complete original proof of it is delivered with it (source0.tex lines 796--944, one \texttt{proof} environment of 149 lines). No mathematics is written, repaired or re-derived: the statement and the proof are the article's own lines, in the article's own notation, and no retained line is substituted or re-flowed.  Retained uncounted as the source-local context the proof needs: lines 329--339; lines 415--418; lines 542--545; lines 717--721.  Nothing was omitted from any retained span, so the package carries no omission record.  No retained line contains a citation, so the wrapper carries no bibliography and no bibliography entry.  No retained line contains a figure environment, an image-inclusion command or drawing code, so no image is delivered and the package declares no asset dependency.  Editorial formatting only, in addition: an AMS \texttt{amsart} preamble that restates the article's own declarations and macros used by the retained lines, namely the environments \texttt{proposition} on separate counters, together with the standard packages those lines need.  The retained environments print in this document as Proposition 1 in order of appearance, whereas the article prints the same items with the numbering of its own full document.  The complete native source of the article as distributed in the arXiv e-print for this exact version is bundled unchanged as \texttt{sources/165979/source0.tex}.
\begin{proposition}
\label{actionofqaut}
Let $\Lambda=(V,E,r,s)$ be a finite directed graph without loops or multiple edges or sources and $C^{\ast}(\Lambda)$ be the corresponding graph $C^{\ast}$ algebra. Then the quantum automorphism group $Aut^{+}(\Lambda)$ acts on $C^{\ast}(\Lambda)$. The action $\alpha:C^{\ast}(\Lambda)\raro C^{\ast}(\Lambda)\otimes C(Aut^{+}(\Lambda))$ is given by, 
\begin{equation*}
   \alpha(p_{i})=\sum\limits_{k=1}^{n}p_{k}\otimes q_{ki}\;\;\;  1\leq i\leq n
\end{equation*}
and 
\begin{equation*}
   \alpha(S_{e_{j}})=\sum\limits_{l=1}^{m}S_{e_{l}}\otimes q_{r(e_{l})r(e_{j})}q_{s(e_{l})s(e_{j})} \;\;\; 1\leq j\leq m.
\end{equation*}
\end{proposition}

\begin{equation}
\label{rep-2}
    (\pi(S_{\lambda}))^{\ast}(\chi_{[\eta]})=\rho(\Lambda)^{-d(\lambda)/2}\chi_{[s(\lambda)]}
\end{equation}

     \begin{equation}
\label{Uformula k lenghth path}
    U(\chi_{[\lambda]})=\sum_{\eta}\chi_{[\eta]}\otimes q_{r(\eta_{1})r(\lambda_{1})}q_{s(\eta_{1})s(\lambda_{1})}\cdots q_{r(\eta_{k})r(\lambda_{k})}q_{s(\eta_{k})s(\lambda_{k})}
\end{equation}

\begin{proposition}
\label{Udensity}
    For the inner product preserving map $U:\mathcal{H}\raro \mathcal{H} \overline{\otimes} C(\text{Aut}^{+}(\Lambda))$,
${\rm Sp}\{U(\xi).q\;|\; \xi\in \mathcal{H}, q\in C(\text{Aut}^{+}(\Lambda))\}\; \text{is dense in}\; \mathcal{H}\overline{\otimes} C(\text{Aut}^{+}(\Lambda))$.
\end{proposition}

\begin{proposition}
    \label{implementation}
    The action $\alpha$ of ${\rm Aut}^{+}(\Lambda)$ on $C^{\ast}(\Lambda)$ as given in Proposition \ref{actionofqaut} is implemented on the spectral data $(C^{\ast}(\Lambda),\clh,D)$ by the unitary co-representation $\widetilde{U}$.
\end{proposition}

\begin{proof}
We start by showing that for any $\lambda,\eta\in\Lambda^{\ast}$, 
\begin{equation}
\label{implementation_equation}
   (\pi\otimes [.]) \alpha(S_{\lambda}^{\ast})U(\chi_{[\eta]})=U(\pi(S_{\lambda}^{\ast})\chi_{[\eta]})
\end{equation}
 \textbf{Case1:} Consider the case $d(\lambda)\geq d(\eta)$. Then there are two subcases: Firstly let $\lambda=\eta\beta$ for some path $\beta$ of finite length. Let's write $\eta=\eta_{1}\eta_{2}\ldots\eta_{m}$ and $\lambda=\lambda_{1}\lambda_{2}\cdots\lambda_{n}$. Then by the unique factorization of paths, $\lambda$ can be written as $\eta_{1}\eta_{2}\ldots\eta_{m}\lambda_{m+1}\ldots\lambda_{n}$ where $\eta_{i}$ and $\lambda_{i}$'s are appropriate composable edges. Note that $s(\lambda)=s(\lambda_{n})$. Then by \ref{rep-2},  
\begin{displaymath}
U(\pi(S_{\lambda}^{\ast})\chi_{[\eta]})=\rho(\Lambda)^{-d(\lambda)/2}U(\chi_{[s(\lambda)]})=\rho(\Lambda)^{-d(\lambda)/2}\sum\limits_{i\in V}\chi_{[i]}\otimes q_{is(\lambda)}\end{displaymath}
 By the Formula \ref{Uformula k lenghth path},
\begin{displaymath}
U(\chi_{[\eta]})=\left(\sum\limits_{\zeta:\zeta=\zeta_{1}\zeta_{2}\ldots\zeta_{m}}\chi_{[\zeta]}\otimes q_{r(\zeta_{1})r(\eta_{1})}q_{s(\zeta_{1})s(\eta_{1})}\cdots q_{r(\zeta_{m})r(\eta_{m})}q_{s(\zeta_{m})s(\eta_{m})}\right)
\end{displaymath}
Therefore, 
\begin{equation*}
    \begin{aligned}
        &(\pi\otimes [.]) \alpha(S_{\lambda}^{\ast})U(\chi_{[\eta]})=  (\pi\otimes [.])(\alpha(S_{\lambda_{n}})^{\ast}\cdots \alpha(S_{\lambda_{1}})^{\ast}) U(\chi_{[\eta]})\\
&=(\pi\otimes [.])\left(\left(\sum\limits_{\xi_{n}}S_{\xi_{n}}^{*}\otimes q_{s(\xi_{n})s(\lambda_{n})}q_{r(\xi_{n})r(\lambda_{n})}\right)\cdots \left(\sum\limits_{\xi_{1}}S_{\xi_{1}}^{\ast}\otimes q_{s(\xi_{1})s(\lambda_{1})}q_{r(\xi_{1})r(\lambda_{1})}\right)\right)\\
&\;\;\;\;\left(\sum\limits_{\zeta}\chi_{[\zeta]}\otimes q_{r(\zeta_{1})r(\eta_{1})}q_{s(\zeta_{1})s(\eta_{1})}\cdots q_{r(\zeta_{m})r(\eta_{m})}q_{s(\zeta_{m})s(\eta_{m})}\right)\\
&=(\pi\otimes [.])\left(\sum\limits_{\substack{\xi=\xi_{1}\cdots\xi_{n}}}S_{\xi}^{*}\otimes q_{s(\xi_{n})s(\lambda_{n})}q_{r(\xi_{n})r(\lambda_{n})}\cdots  q_{s(\xi_{1})s(\lambda_{1})}q_{r(\xi_{1})r(\lambda_{1})} \right)\\
&\;\;\;\;\left(\sum\limits_{\zeta}\chi_{[\zeta]}\otimes q_{r(\zeta_{1})r(\eta_{1})}q_{s(\zeta_{1})s(\eta_{1})}\cdots q_{r(\zeta_{m})r(\eta_{m})}q_{s(\zeta_{m})s(\eta_{m})}\right)
    \end{aligned}
\end{equation*}
\begin{equation*}\hfill
    \begin{aligned}
&=\left(\sum\limits_{\substack{\xi=\xi_{1}\cdots\xi_{n}}}\pi(S_{\xi})^{\ast}\otimes [q_{s(\xi_{n})s(\lambda_{n})}q_{r(\xi_{n})r(\lambda_{n})}\cdots  q_{s(\xi_{1})s(\lambda_{1})}q_{r(\xi_{1})r(\lambda_{1})}] \right)\\
&\;\;\;\;\left(\sum\limits_{\zeta}\chi_{[\zeta]}\otimes q_{r(\zeta_{1})r(\eta_{1})}q_{s(\zeta_{1})s(\eta_{1})}\cdots q_{r(\zeta_{m})r(\eta_{m})}q_{s(\zeta_{m})s(\eta_{m})}\right)
\\
\end{aligned}
\end{equation*} 
\begin{equation*}\hfill
    =\sum\limits_{\substack{\xi,\zeta\\ d(\xi)\geq d(\zeta)}}\pi(S_{\xi})^{*}\chi_{[\zeta]}\otimes 
\begin{aligned}[t]
&q_{s(\xi_{n})s(\lambda_{n})}q_{r(\xi_{n})r(\lambda_{n})}\cdots  q_{s(\xi_{1})s(\lambda_{1})}q_{r(\xi_{1})r(\lambda_{1})}q_{r(\zeta_{1})r(\eta_{1})}\\
&q_{s(\zeta_{1})s(\eta_{1})}\cdots q_{r(\zeta_{m})r(\eta_{m})}q_{s(\zeta_{m})s(\eta_{m})}
\end{aligned}
\end{equation*}
Note that $d(\xi)\geq d(\zeta)$ as the action $\alpha$ and the map $U$ are both degree preserving. Now for any path $\xi,\zeta$ with $d(\xi)\geq d(\zeta)$ such that $\xi=\zeta\gamma$ for some path $\gamma$, $\pi(S_{\xi}^{\ast})(\chi_{[\zeta]})=\chi_{[s(\xi)]}$. If there is no path $\gamma$ of length $(n-m)$ such that $\xi=\zeta\gamma$, then $\pi(S_{\xi}^{\ast})(\chi_{[\zeta]})=0$. For $\xi=\zeta\gamma$, using the unique factorisation of paths, we write $\xi=\zeta_{1}\zeta_{2}\cdots \zeta_{m}\xi_{m+1}\cdots \xi_{n}$, $\gamma=\xi_{m+1}\xi_{m+2}\cdots \xi_{n}$. Note that $s(\xi)=s(\xi_{n})$. Therefore, using the expression $\lambda=\eta_{1}\cdots \eta_{m}\lambda_{m+1}\cdots\lambda_{n}$, the last summation becomes
\begin{equation*}
    \rho(\Lambda)^{-d(\xi)/2}\sum\limits_{\substack{\xi,\zeta\\ d(\xi)\geq d(\zeta) \\ \xi=\zeta\gamma}}\chi_{[s(\xi)]}\otimes 
\begin{aligned}[t]
&q_{s(\xi_{n})s(\lambda_{n})}q_{r(\xi_{n})r(\lambda_{n})}
\cdots q_{s(\xi_{m+1})s(\lambda_{m+1})}q_{r(\xi_{m+1})r(\lambda_{m+1})}\\
&q_{s(\zeta_{m})s(\eta_{m})}q_{r(\zeta_{m})r(\eta_{m})}
\cdots
q_{s(\zeta_{1})s(\eta_{1})}q_{r(\zeta_{1})r(\eta_{1})}\\
&q_{r(\zeta_{1})r(\eta_{1})}q_{s(\zeta_{1})s(\eta_{1})}\cdots\cdots q_{r(\zeta_{m})r(\eta_{m})}q_{s(\zeta_{m})s(\eta_{m}).}
\end{aligned}
\end{equation*}
As $d(\xi)=d(\lambda)$, ignoring the common scalar factor $\rho(\Lambda)^{-d/2}$, we shall show that for a fixed $i\in V$, the coefficient of $\chi_{[i]}$ in the above summation is $q_{is(\lambda)}$ proving \ref{implementation_equation} in this subcase. The coefficient is given by
\begin{equation*}
    \sum\limits_{\substack{\xi,\zeta\\ n=d(\xi)\geq d(\zeta)=m \\ \xi=\zeta\gamma, s(\xi)=i}}
\begin{aligned}[t]
&q_{s(\xi_{n})s(\lambda_{n})}q_{r(\xi_{n})r(\lambda_{n})}
\cdots q_{s(\xi_{m+1})s(\lambda_{m+1})}q_{r(\xi_{m+1})r(\lambda_{m+1})}q_{s(\zeta_{m})s(\eta_{m})}q_{r(\zeta_{m})r(\eta_{m})}
\\
&\cdots q_{s(\zeta_{1})s(\eta_{1})}q_{r(\zeta_{1})r(\eta_{1})}q_{r(\zeta_{1})r(\eta_{1})}q_{s(\zeta_{1})s(\eta_{1})}\cdots q_{r(\zeta_{m})r(\eta_{m})}q_{s(\zeta_{m})s(\eta_{m})}
\end{aligned}
\end{equation*}
Now consider the term $q_{s(\zeta_{1})s(\eta_{1})}q_{r(\zeta_{1})r(\eta_{1})}q_{s(\zeta_{1})s(\eta_{1})}$. As $\zeta_{1}$ varies over all edges such that $\zeta_{1}\zeta_{2}\ldots\zeta_{m}$ is an $m$-length path, as in the proof of \ref{Udensity}, we can replace this term in the summation with $\sum_{k,l\in V}q_{ks(\eta_{1})}q_{lr(\eta_{1})}q_{ks(\eta_{1})}$ which is $1$ as $\sum_{k}q_{ki}=1$ for any $i$. Therefore, the summation reduces to
\begin{equation*}
    \sum\limits_{\substack{\xi,\zeta\\ d(\xi)\geq d(\zeta) \\
s(\xi)=i}}
\begin{aligned}[t]
&\;q_{s(\xi_{n})s(\lambda_{n})}q_{r(\xi_{n})r(\lambda_{n})}\cdots
q_{s(\xi_{m+1})s(\lambda_{m+1})}q_{r(\xi_{m+1})r(\lambda_{m+1})}
\\
&\;q_{s(\zeta_{m})s(\eta_{m})}q_{r(\zeta_{m})r(\eta_{m})}\cdots q_{s(\zeta_{2})s(\eta_{2})}q_{r(\zeta_{2})r(\eta_{2})}q_{s(\zeta_{2})s(\eta_{2})}\\
&\;\cdots q_{r(\zeta_{m})r(\eta_{m})}q_{s(\zeta_{m})s(\eta_{m})}
\end{aligned}
\end{equation*}
By using the same trick repeatedly, the above summation can further be reduced to
\begin{equation*}
    \sum\limits_{\substack{\xi,
s(\xi)=i}}q_{s(\xi_{n})s(\lambda_{n})}q_{r(\xi_{n})r(\lambda_{n})}\cdots q_{s(\xi_{m+1})s(\lambda_{m+1})}q_{r(\xi_{m+1})r(\lambda_{m+1})}
\end{equation*}
Now note that in the above expression, $(r(\xi_{n}),s(\xi_{n-1}),\cdots, s(\xi_{m+1}),r(\xi_{m+1}))$ can be replaced by $(l_{n},k_{n-1},l_{n-1},\cdots,k_{m+1},l_{m+1})$ where $l_{i}'s$ and $k_{i}'s$ vary over the vertex set. This is because if $((k_{m+1},l_{m+1}),\cdots,(k_{n-1},l_{n-1}),(i,l_{n}))$ forms a path of degree $n-m$ starting at $i$ then the corresponding coefficients already appear in the summation. Otherwise, as $\lambda_{m+1}\cdots\lambda_{n}$ is an $(n-m)$-length path, \begin{displaymath}q_{is(\lambda_{n})}q_{l_{n}r(\lambda_{n})}\cdots q_{l_{m+1}r(\lambda_{m+1})}=0.\end{displaymath} Therefore, the coefficient of $\chi_{[i]}$ in the summation becomes
\begin{equation*}
    \sum\limits_{\substack{l_{n},k_{n-1},\ldots,k_{m+1},l_{m+1}\\s(\xi)=s(\xi_{n})=i}}q_{s(\xi_{n})s(\lambda_{n})}q_{l_{n}r(\lambda_{n})}\cdots q_{k_{m+1}s(\lambda_{m+1})}q_{l_{m+1}r(\lambda_{m+1})}=q_{is(\lambda)}.
\end{equation*}
Now consider the sub-case where there is no path $\beta$ such that $\lambda=\eta\beta$ which means there exists $k\in\{1,2,\cdots,m\}$ such that $\lambda_{k}\neq \eta_{k}$  and hence $\pi(S_{\lambda}^{\ast})\chi_{[\eta]}=0$. For simplicity let's assume that $\lambda_{1}\neq \eta_{1}$. The case $k>1$ can be proved along the same line of argument. Using the same computation as above the term $(\pi\otimes [.]) \alpha(S_{\lambda}^{\ast})U(\chi_{[\eta]})$ in the LHS is expanded as,
\begin{equation*}
    \rho(\Lambda)^{-d/2}\sum\limits_{\substack{\xi,\zeta\\ d(\xi)\geq d(\zeta) \\ \xi=\zeta\gamma}}\chi_{[s(\xi)]}\otimes 
\begin{aligned}[t]
&q_{s(\xi_{n})s(\lambda_{n})}q_{r(\xi_{n})r(\lambda_{n})}
\cdots q_{s(\xi_{m+1})s(\lambda_{m+1})}q_{r(\xi_{m+1})r(\lambda_{m+1})}\\
&q_{s(\zeta_{m})s(\lambda_{m})}q_{r(\zeta_{m})r(\lambda_{m})}
\cdots
q_{s(\zeta_{1})s(\lambda_{1})}q_{r(\zeta_{1})r(\lambda_{1})}q_{r(\zeta_{1})r(\eta_{1})}\\
&q_{s(\zeta_{1})s(\eta_{1})}\cdots\cdots q_{r(\zeta_{m})r(\eta_{m})}q_{s(\zeta_{m})s(\eta_{m})}.
\end{aligned}
\end{equation*}
As 
$\lambda_{1}\neq \eta_{1}$ that means either $r(\lambda_{1})\neq r(\eta_{1})$ or $s(\lambda_{1})\neq s(\eta_{1})$. Now for the case $r(\lambda_{1})\neq r(\eta_{1})$ the term $q_{r(\zeta_{1})r(\lambda_{1})}q_{r(\zeta_{1})r(\eta_{1})}=0$ otherwise lets assume that $r(\lambda_{1})= r(\eta_{1})$ and $s(\lambda_{1})\neq s(\eta_{1})$, 
now using the same argument as before, we can reduce the term $q_{s(\zeta_{1})s(\lambda_{1})}q_{r(\zeta_{1})r(\eta_{1})}q_{s(\zeta_{1})s(\eta_{1})}$ to $\sum\limits_{k,l}q_{ks(\lambda_{1})}q_{lr(\eta_{1})}q_{ks(\eta_{1})}=\sum\limits_{k}q_{ks(\lambda_{1})}q_{ks(\eta_{1})}=0$
and we are done.\\
\textbf{Case 2:} $d(\lambda)<d(\eta)$: Again as in case 1, it can be divided into two subcases: there is some $\beta$ such that $\lambda\beta=\eta$ and there is no $\beta$ such that $\lambda\beta=\eta$. In the first subcase, $\pi(S_{\lambda}^{\ast})(\chi_{[\eta]})=\chi_{[\beta]}$. In the second subcase $\pi(S_{\lambda}^{\ast})(\chi_{[\eta]})=0$. In both the cases, using the same computational technique, it can be shown that 
\begin{displaymath}
    U(\pi(S_{\lambda}^{\ast})(\chi_{[\eta]}))=(\pi\otimes[.])\alpha(S_{\lambda}^{\ast})U(\chi_{[\eta]}).
\end{displaymath}

Therefore, for any path $\eta\in\Lambda^{\ast}$, we have $(\pi\otimes[.])\alpha(S_{\lambda}^{\ast})U(\chi_{[\eta]})=U(\pi(S_{\lambda^{\ast}}))(\chi_{[\eta]})$ and hence for any $q\in C({\rm Aut}^{+}(\Lambda))$,
\begin{displaymath}
  (\pi\otimes[.])\alpha(S_{\lambda}^{\ast})\widetilde{U}(\chi_{[\eta]}\otimes q)=\widetilde{U}(\pi(S_{\lambda}^{\ast})\otimes [1])(\chi_{\eta}\otimes q). 
\end{displaymath} 
By density of ${\rm Sp}\{\chi_{[\eta]}\otimes q:\eta\in\Lambda^{\ast}, q\in C({\rm Aut}^{+}(\Lambda))\}$ in the Hilbert space $\clh\otimes L^{2}({\rm Aut}^{+}(\Lambda))$ and unitarity of $\widetilde{U}$, we have for any $\lambda\in\Lambda^{\ast}$,
\begin{eqnarray}
\label{implementation star}
    (\pi\otimes[.])\alpha(S_{\lambda}^{\ast})=\widetilde{U}(\pi(S_{\lambda}^{\ast})\otimes[1])\widetilde{U}^{\ast}\in \clb(\clh\otimes L^{2}({\rm Aut}^{+}(\Lambda))).
\end{eqnarray}
Using the same computational technique it can be shown that 
\begin{eqnarray}
\label{implementation non star}
    (\pi\otimes[.])\alpha(S_{\lambda})=\widetilde{U}(\pi(S_{\lambda})\otimes[1])\widetilde{U}^{\ast}\in \clb(\clh\otimes L^{2}({\rm Aut}^{+}(\Lambda)))
\end{eqnarray}
for any $\lambda\in\Lambda^{\ast}$. Alternatively, one can see as before that Equation \ref{implementation non star} is equivalent to the following for $\lambda,\eta\in\Lambda^{\ast}$:
\begin{equation}
\label{non star}
    (\pi\otimes[.])\alpha(S_{\lambda})U(\chi_{[\eta]})=U(\pi(S_{\lambda}))(\chi_{[\eta]}).
\end{equation}
Again recalling the representation $\pi$ of $C^{\ast}(\Lambda)$, we see that RHS of the above is equal to $\rho(\Lambda)^{d(\lambda/2)}U(\chi_{[\lambda\eta]})$ whereas the LHS is equal to 
\begin{equation}
\label{temp}    
\begin{aligned}
    &\rho(\Lambda)^{d(\mu)/2}\sum_{}\chi_{[\mu\xi]}\otimes q_{r(\mu_{1})r(\lambda_{1})}q_{s(\mu_{1})s(\lambda_{1})}\cdots q_{r(\mu_{n})r(\lambda_{n})}q_{s(\mu_{n})s(\lambda_{n})}q_{r(\xi_{1})r(\eta_{1})}\\
    &q_{s(\xi_{1})s(\eta_{1})}\cdots 
    q_{r(\xi_{m})r(\eta_{m})}q_{s(\xi_{m})s(\eta_{m})}
\end{aligned}
\end{equation} 
where $\mu=\mu_{1}\ldots\mu_{n}$ and $\xi=\xi_{1}\ldots\xi_{m}$;
\begin{displaymath}  
\alpha(S_{\lambda})=\sum S_{\mu}\otimes q_{r(\mu_{1})r(\lambda_{1})}q_{s(\mu_{1})s(\lambda_{1})}\cdots q_{r(\mu_{n})r(\lambda_{n})}q_{s(\mu_{n})s(\lambda_{n})}
\end{displaymath} 
and 
\begin{displaymath}
U(\chi_{[\eta]})=\sum\limits_{\xi}\chi_{[\xi]}\otimes q_{r(\xi_{1})r(\eta_{1})}q_{s(\xi_{1})s(\eta_{1})}\cdots \\
    q_{r(\xi_{m})r(\eta_{m})}q_{s(\xi_{m})s(\eta_{m})}.
\end{displaymath} 
A moment's reflection reveals that the coefficients of $\chi_{[\mu\xi]}$ in \ref{temp} are the coefficients of $S_{\mu}S_{\xi}$ in the expression of $\alpha(S_{\lambda})\alpha(S_{\eta})$. Then the homomorphic property of $\alpha$ coupled with the fact that $d(\mu)=d(\lambda)$ for all $\mu$ proves Equation \ref{non star}.  
Since both $(\pi\otimes[.])\alpha$ and ${\rm ad}_{\widetilde{U}}$ are $C^{\ast}$-homomorphisms, we have the following from \ref{implementation star} and \ref{implementation non star}:
\begin{displaymath}
 (\pi\otimes[.])\alpha(S_{\lambda}S_{\mu}^{\ast})=\widetilde{U}(\pi(S_{\lambda}S_{\mu}^{\ast})\otimes[1])\widetilde{U}^{\ast}\in \clb(\clh\otimes L^{2}({\rm Aut}^{+}(\Lambda))),   
\end{displaymath}
for $\lambda,\mu\in\Lambda^{\ast}$. As ${\rm Sp}\{S_{\lambda}S_{\mu}^{\ast}:\lambda,\mu\in\Lambda^{\ast}\}$ is dense in $C^{\ast}(\Lambda)$, using the continuity of the maps involved, we get
\begin{displaymath}
 (\pi\otimes[.])\alpha(a)=\widetilde{U}(\pi(a)\otimes[1])\widetilde{U}^{\ast}\in \clb(\clh\otimes L^{2}({\rm Aut}^{+}(\Lambda))),   
\end{displaymath}
for all $a\in C^{\ast}(\Lambda)$, proving the proposition.
\end{proof} 

\end{document}
